\chapter{绪论：为什么研究一只捧腹大笑的青蛙}\label{ch:01}

\begin{nlkeypoints}
  \begin{enumerate}[label=(\arabic*),leftmargin=2.4em]
    \item “奶蛙”不只是一个网络热梗，而是具备可考据的来源线索、有规模的二创生态、
          有产业接口与跨校扩散路径的\keyw{文化现象}，足以作为研究对象进入严肃讨论。
    \item 本书把“官方奶龙”“奶蛙”“捧腹大笑”区分为三个不同层面的对象，
          并在全书中坚持\keyw{证据分级}（A 一手权威 / B 机构公开 / C 社群考据 / M 本书自建模型），
          凡属推测一律显式标注。
    \item 本书回答四个递进的研究问题：它是什么符号系统？人们为什么笑？
          它如何扩散？这类现象应当如何研究？
  \end{enumerate}
\end{nlkeypoints}

2026 年，一只腹部膨大、四肢细小、黄皮绿眼的生物，以一种异常亢奋的笑声
占领了中文互联网的相当一部分信息流。它既不是任何一部动画里的角色，
也不是任何一个账号的“主人”，却拥有完整的形象谱系、稳定的“台词”、
成体系的动作、遍布高校校园的实体立牌，甚至出现了未经授权的周边商品。
有二手评述称其为“2026 年传播最广、影响最深、二创最活跃的中文互联网
meme”\conf{C}\cite{shoutanji2026naiwa}。本书把这一现象称为\term{奶蛙}，
把构成其核心表达的笑声称为\term{捧腹大笑}。

本章要做三件事：描述这个现象；说明它为何构成一个值得追问的“谜题”；
交代本书的研究问题、结构与方法承诺。

\section{一只笑了一整年的青蛙：现象描述}\label{sec:01:phenomenon}

先看材料本身。奶蛙的典型形态是：一只近似青蛙的生物，通体黄色，腹部显著膨大，
四肢短小，双手呈黑色；它几乎总是处于大笑状态，笑声经过变调处理，
异常响亮且带有某种“绷不住”的质感\conf{C}\cite{shoutanji2026naiwa}。
在各大平台的二创视频中，这只生物频繁出现在各种“抽象”桥段的结尾，
用一次集体大笑为视频收尾——这种“结尾一起笑”的固定套路本身，
后来也成了该迷因的形式特征之一\conf{C}\cite{shoutanji2026naiwa}。

与现象同步存在的，是它的“正典”背景。奶蛙形象与官方动画 IP
\term{奶龙}高度相似：奶龙由第七印象文化传媒（深圳）有限公司运营，
2017 年团队创立，2019 年构思形象，2020 年 1 月 22 日上线短视频平台
\conf{A}\cite{jiang2026nailong}\conf{B}；截至 2026 年初，其全网粉丝超 3800 万，
授权合作品牌超 200 家，衍生产品 SKU 逾 2000 个，产品销往 57 个国家和地区
\conf{A}\cite{jiang2026nailong}。也就是说，一只“不属于任何官方”的青蛙，
寄生在一个极其成功、且有专业法务团队的正版 IP 旁边生长——
这个张力将在第九章详细展开。

\begin{nlhistory}[史料与考据]
  关于奶蛙的形象来源，目前最系统的社群考据指向：2024 年 10 月 2 日，
  星野 AI 平台用户“纲手【百豪】”创建“奶龙”智能体时，平台 AI 生成的头像
  出现了体型畸变，形成最初版本奶蛙形象；2026 年 8 月前后，
  B 站用户“雾蚀心忘症”“看仙碎”等爱好者完成了对该智能体
  （ID \fname{zNCvQDKfXy}）的确认，但该智能体现已不可访问\conf{C}
  \cite{lostmedia2026naiwa}。注意：这是\keyw{考据}，不是定论——
  考据词条自陈“仅为推测、缺乏直接证据”。第四章将以整章篇幅
  示范如何把一条“未结案的证据链”当作教学材料。
\end{nlhistory}

需要强调的第一点：\textbf{奶蛙不是“一个梗”，而是三个层面的叠加}。
第一层是官方奶龙——一部定位“治愈、欢乐”的子供向动画及其完整 IP 体系
\conf{A}\cite{jiang2026nailong}；第二层是民间的奶蛙形象——一个 AI 畸变产物
经过社群筛选后稳定下来的“变种”；第三层是“捧腹大笑”——
把笑声、形象、动作与特定情境捆绑之后形成的完整表达形式。
全书将始终把这三层区分开，混用它们会产生大量错误解释
（比如把官方 IP 的成功直接归因于奶蛙，或把奶蛙的传播能力
当作官方奶龙的传播能力）。

\section{一个看似违背常理的谜题}\label{sec:01:puzzle}

第一版教材的绪论提出过一个问题：“为什么越聪明的人反而越喜欢那些
无厘头、不着边际的东西？”这个问题值得保留，因为它指出了本现象
最耐人寻味的特征：奶蛙的高频使用者与高频二创者，高度集中于高校环境，
尤其是竞争压力大的顶尖院校——同好会、立牌、校园表情包
（“数分奶蛙”“捧着航概的奶蛙”）都生长在校园空间里\conf{C}。
如果把“高智识人群 × 无厘头内容”当成一个待解释的组合，
它至少面临三种竞争读法：

\begin{enumerate}[label=\arabic*.]
  \item \textbf{纯消遣读法}：它只是一种低成本、低门槛的娱乐，
        大学生恰好有空刷手机，与“聪明”无关，是幸存者偏差。
  \item \textbf{技术—平台读法}：扩散主要由平台算法与 AIGC 工具推动，
        个体的“喜欢”是被动卷入，与心理机制无关。
  \item \textbf{功能—符号读法}：捧腹大笑承担某种系统性的情绪与
        社交功能（代偿、稀释、表态），高校恰好是这种功能
        需求最密集、传播介质最便利的场域。
\end{enumerate}

本书采取第三种读法作为主解释，但\keyw{全程保留前两种作为竞争性解释}：
第七章会给出每个机制的可检验假设与证伪条件，第十一章给出问卷设计。
“功能—符号读法”并不是本书的结论，而是本书\keyw{要检验的立场}。
在检验完成之前，它只是工作假设。这一点是本书与大多数“梗评”
文章的根本区别，也是本书的方法论承诺所在。

为什么功能读法值得认真对待？因为该现象的几个特征难以用纯消遣解释：
其一，它的使用有明显的\textbf{情境规律}——它几乎总是出现在“压力桥段”
之后，而不是随机位置；其二，它伴随稳定的\textbf{自我呈现策略}——
用一张梗图代替一句自述，把“我很难受”翻译成“我们都在笑”；
这与戈夫曼所述前台/后台管理中的“印象缓冲”结构高度同构
\cite{goffman1959presentation}；其三，它的表达门槛极低，
而在线环境的匿名与去抑制特征进一步降低了使用成本\cite{suler2004disinhibition}。
这些观察都还只是现象学层面的，机制分析见第七章。

\begin{nltermbox}{术语：meme（迷因）}
  \en{meme} 一词由道金斯在《自私的基因》第 11 章提出，
  指文化信息的复制单位，与基因类比：可复制、可变异、受选择压力
  \cite{dawkins1976selfish}。数字时代的分析以 Shifman 的三维框架
  （内容 / 形式 / 立场）为代表\cite{shifman2013memes}。
  本书在第五章给出该概念的操作化使用方式，第六章用它解剖奶蛙。
\end{nltermbox}

\section{本书的四个研究问题}\label{sec:01:questions}

全书围绕四个递进问题展开：

\begin{enumerate}[label=\textbf{Q\arabic*},leftmargin=4em,itemsep=0.4em]
  \item \textbf{符号问题}：奶蛙是一个怎样的符号系统？
        它的内容、形式与立场分别是什么？（第六章，以 Shifman
        三维框架\cite{shifman2013memes}为工具）
  \item \textbf{机制问题}：为什么大学生会笑，而且愿意反复转发？
        五种候选机制（情绪代偿、群体稀释、反精致、校园社交、
        逆反认证）各自能解释什么、不能解释什么？（第五章、第七章）
  \item \textbf{扩散问题}：它如何从短视频平台扩散到高校校园与商品世界？
        正典奶龙的产业史、奶蛙的起源考据、AIGC 生产链、
        周边经济与校园空间实践分别贡献了哪一段？（第三、四、八、九、十章）
  \item \textbf{方法问题}：研究一个“活着变”的网络文化现象，
        证据应当如何分级，考据应当如何伦理化，结论应当如何标注
        脆弱性？（第五章、第十二章）
\end{enumerate}

四个问题不是并列的：Q1 是 Q2 的输入（不懂符号系统就谈不上机制），
Q2 是 Q3 的解释层（扩散需要动力机制），Q4 贯穿全程并收束于
第十二章。这个依赖结构决定了章节顺序。

\section{本书的结构}\label{sec:01:structure}

全书十二章分三部分。第一部分（第 2—5 章）确立对象与工具：
界定“正典—变种”的对象结构，梳理正典奶龙的产业史与奶蛙的起源考据，
给出迷因理论框架与本书自建的传播教学模型。第二部分（第 6—9 章）
进入核心分析：解剖符号系统，把“为什么笑”重写为可检验的假设，
再讨论 AIGC 生产链与文化经济如何把“玩梗”接上产业。
第三部分（第 10—12 章）落到实践与方法：高校亚文化的空间实践与组织形态、
受体研究的问卷与访谈设计、以及全书方法论的收束（证据分级、考据伦理、
版权、可复现性）。

\begin{table}[htbp]\centering
  \caption{全书三部分与十二章的功能分工}\label{tab:01:structure}
  \small
  \begin{tabularx}{\textwidth}{@{}llXX@{}}
    \toprule
    部分 & 章 & 回答的问题 & 主要材料 \\ \midrule
    一：对象与工具
      & 第 2 章 & 奶龙学研究的对象与边界是什么 & 概念界定 \\
      & 第 3 章 & 正典奶龙如何成为一个产业 & 媒体公开报道（A/B 级） \\
      & 第 4 章 & 奶蛙从何而来（未结案的证据链） & 社群考据（C 级） \\
      & 第 5 章 & 用哪些理论工具分析迷因 & 理论文献（A 级） \\ \midrule
    二：符号与机制
      & 第 6 章 & 奶蛙是什么符号系统 & 形象—语言—动作语料 \\
      & 第 7 章 & 为什么笑：机制与假设 & 机制矩阵 + 假设 \\
      & 第 8 章 & AIGC 如何降低生产成本 & 生产链拆解 \\
      & 第 9 章 & 玩梗如何接入产业 & 授权与二创周边 \\ \midrule
    三：实践与方法
      & 第 10 章 & 奶蛙如何进入校园空间与组织 & 高校案例 \\
      & 第 11 章 & 大学生为什么笑：受体研究设计 & 问卷与访谈设计 \\
      & 第 12 章 & 这类研究应当怎么做 & 方法论与伦理 \\ \bottomrule
  \end{tabularx}
\end{table}

\section{方法承诺：事实、说法与模型}\label{sec:01:method}

本书涉及三类陈述，全书用统一标识区分：

\begin{enumerate}[label=\arabic*.]
  \item \textbf{事实}（\conf{A} / \conf{B}）：可追溯一手权威来源或
        机构公开信息的陈述。引用格式遵循 GB/T 7714—2015\cite{gbt7714_2015}。
  \item \textbf{说法}（\conf{C}）：社群考据、二手评述。只能以
        “据……”“有说法认为”引入，且不得作为论证前提。
  \item \textbf{模型}（\conf{M}）：本书为教学目的自建的框架、
        方程与示例数据。必须显式声明“本书构造”“非经验事实”，
        且参数取值说明为示意值。
\end{enumerate}

这一分级制度是本书最重要的体例装置，其设计理由、操作细则与局限
在第十二章系统讨论。读者若发现某条 \conf{C} 级陈述在正文中
被当作事实使用，那正是本书鼓励你指出的错误。

\begin{nlcritique}[批判性反思]
  研究“乐子”有一个天然的两难：端着架子，分析会失去对象；
  跟着一起乐，分析会失去距离。本书的立场是\keyw{对现象保持幽默，
  对方法保持严肃}——章名可以玩笑，表格必须严谨；脚注可以吐槽，
  结论必须可核查。同时必须声明：奶蛙二创中存在未经授权的形象再生产，
  这在版权法上是有争议的（第九章讨论）；本书把这一现象当作
  研究对象，不构成对任何侵权行为的背书。
\end{nlcritique}

\section{本章小结}\label{sec:01:summary}

本章完成了三件事。第一，把研究对象拆成“官方奶龙—民间奶蛙—
捧腹大笑”三个层面，这是全书一切分析的单位约定。第二，
把第一版教材“为什么聪明人爱看无厘头”的问题重写为
三种竞争读法，并明确本书的工作立场：以功能—符号读法为主解释，
以纯消遣与技术—平台读法为竞争性解释，检验工具在第七、十一章。
第三，宣布全书的证据分级制度（A/B/C/M）与四个递进研究问题。

尚未解决的问题：奶蛙的符号结构（第六章）、五种机制的
可检验化（第七章）、以及一个更深的认识论问题——
当研究对象本身以“解构严肃性”为乐时，一本严肃教材
如何保持自身立场的一致性。这个问题在第十二章收束。

\begin{nlexercises}
  \item 选取你最近 30 天内实际转发或使用过的一个梗，
        用“内容 / 形式 / 立场”三维各写一句话描述，
        并指出其中哪一维是社群硬约定（用错会被识别为外行），
        哪一维可自由发挥。
  \item 为第七章的“压力源假设”与“平台算法假设”各设计一条
        \keyw{可观察指标}（必须是单一数字或单一频率，能直接计数），
        并说明若两条指标的相关方向相反，说明什么。
  \item 在本书“事实手册”中找到两个来自不同来源、
        数值不同的同一指标（例如粉丝数），解释为什么本书
        选择并列而非取均值，并讨论这种处理对引用者的影响。
\end{nlexercises}

\nlkeywords 奶蛙、奶龙、捧腹大笑、meme（迷因）、证据分级、
参与式文化、竞争性解释
